\documentclass[11pt]{article}
\usepackage[margin=2.5cm]{geometry}
\usepackage{booktabs}
\usepackage{amsmath}
\usepackage{hyperref}
\usepackage{xcolor}
\usepackage{parskip}

\title{Proposal follow-up: Stage 2 (multimodal signals) and Stage 3 (cross-country panel)\\
\large Status report on \emph{``Remittance shock 2022: when both the National Bank and the World Bank got it wrong''}}
\author{}
\date{\today}

\begin{document}
\maketitle

\section*{Where this sits relative to the original proposal}

The original proposal laid out three stages: Stage 1 (baseline models on Kyrgyzstan data),
Stage 2 (multimodal covariates -- were there early leading indicators for the 2022 shock?), and
Stage 3 (cross-country panel with Tajikistan and Armenia, to avoid the ``anecdotal single case''
risk the proposal itself flagged). Stage 1 is done and is the MIDI~2026 submission, already
through one peer-review pass and page-limited to 8pp -- it is not revisited here. This report
covers what was attempted for Stages 2 and 3 after that submission was finished, what was found,
and what remains open. All code and data referenced below are committed to the project
repository (\texttt{capitalisming/midi2026}, branch \texttt{main-dlkex3}).

\section{Stage 2: multimodal signals}

\subsection{What was asked}
The proposal's Stage 2 question was direct: with covariates available in real time in
February--March 2022 (ruble exchange rate, payment-system transaction volumes, border-service
statistics), could the 2022 shock have been anticipated -- or was it genuinely unpredictable
from the information available at the time?

\subsection{What data were used}
Border-service statistics and payment-system transaction data were not obtainable in a reusable
form; the covariates actually available and already downloaded in an earlier phase of the
project were FX (RUB/USD, Central Bank of Russia), Brent crude oil (EIA), and the MOEX equity
index -- proxies for the same underlying shock (the Russia-sanctions channel), not the exact
series the proposal named, but the closest available.

\subsection{The leakage problem, and why it was previously blocked}
A prior scope note in the codebase (\texttt{panel.py}) recorded that covariates had been tried
and set aside: a naively-lagged covariate feature can leak. Concretely, Brent's own publication
lag (about one month) means that at forecast issue date $v$, the freshest Brent figure known is
for month $v{-}1$, not month $v$ -- so a feature defined as ``Brent at the month before the
target'' would, for the query row specifically, reach for a value that does not yet exist.

What this follow-up found: the leakage-safe machinery to handle this correctly
(\texttt{realtime.covariates\_at()}, with a documented lag per covariate and its own dedicated
leakage test) already existed in the codebase and was already passing tests -- it had simply
never been wired into the feature matrix the ridge models train on. The fix
(\texttt{panel.\allowbreak make\_covariate\_features()}) shifts each covariate's log-change feature by
\emph{(1 + that covariate's own publication lag)} months, so the query row never reaches past
what is actually known at $v$. A new black-box leakage test tampers every covariate value dated
after its known-at-$v$ cutoff with an absurd sentinel and confirms the rebuilt row is unchanged
-- i.e. the feature construction provably never reads data it should not have access to. All
1{,}610 tests pass (1{,}553 pre-existing + 57 new).

\subsection{Result}
Same protocol as the rest of the paper: grid search $\alpha$ over the development period only,
freeze, then evaluate.

\begin{table}[h]
\centering
\begin{tabular}{lcc}
\toprule
 & Main window (n=50) & Shock 2022 (n=8) \\
\midrule
Expanding ridge, no covariates & 32.7 & 53.4 \\
Expanding ridge + FX/Brent/IMOEX & 34.7 & 59.9 \\
\bottomrule
\end{tabular}
\caption{MAE (million USD). Covariates were tuned on the development period ($\alpha=0.1$,
dev MAE 26.7) using the same alpha grid as the no-covariate baseline.}
\end{table}

\textbf{The covariates make forecasts worse, not better} -- both in the ordinary period and,
more importantly, during the shock itself. This is a direct, negative answer to the proposal's
Stage 2 question: at least for this set of signals (FX, oil, a Moscow equity index) under a
linear specification, there is no evidence of an early-warning advantage. This is reported as a
finding, not a failure -- it rules out the more optimistic reading of the 2022 episode (``the
signal was there, institutions just missed it'') in favour of the harder one already central to
the main paper: the shock looked, from the available leading indicators, like noise.

\subsection*{Caveats}
\begin{itemize}
\item Only three covariates were available; the proposal's own suggestions (payment-system
  transaction volumes, border-crossing statistics) could behave differently and were not tested.
\item A linear (ridge) specification was used throughout, for consistency with the rest of the
  paper; a model that captures nonlinear covariate interactions might do better (or could
  equally overfit worse -- untested).
\end{itemize}

\section{Stage 3: cross-country panel}

\subsection{Tajikistan: not feasible with this methodology}
The proposal named Tajikistan and Armenia. Tajikistan was investigated first and set aside:
\begin{itemize}
\item The National Bank of Tajikistan's public statistical bulletin archive
  (\texttt{nbt.tj/en/statistics/\allowbreak statistical\_bulletin.php}) keeps only \emph{one issue per year}
  for 2011--2025 -- not the dense monthly captures the real-time protocol needs to reconstruct
  what was known at each historical issue date. Only the current year (2026) is available
  month-by-month.
\item Separately, a July~2026 news item reported that the National Bank of Tajikistan had
  stopped disclosing remittance figures at all, proposing to replace the term with
  ``financial inflows'' without giving numbers -- a live transparency problem, not just an
  archival one.
\item Official balance-of-payments data for Tajikistan is published annually, with results
  ``typically ready for publication by April'' of the following year -- a fundamentally
  different (and far coarser) cadence than the monthly, chronologically-reconstructed protocol
  this project is built around.
\end{itemize}
Conclusion: Tajikistan cannot support the same real-time methodology used for Kyrgyzstan. A
weaker design (year-end snapshots only) was considered and rejected as not worth doing --
it would answer a different, much less interesting question.

\subsection{Armenia: data sourced, first result obtained}

\paragraph{Series used.} The Central Bank of Armenia (CBA) publishes two different
remittance-related series:
\begin{enumerate}
\item An official quarterly balance-of-payments remittances series (IMF BPM6 methodology,
  since 2004), with a \emph{documented} publication lag (``last working day of the third month
  following the reporting period'') and known to be revised. Wrong frequency for this project's
  $h=1$-month design -- quarterly data would give roughly 8 shock-window origins across the
  entire sample, not per shock.
\item A monthly table, \emph{``Total money transfers of individuals sent to and received from
  abroad through commercial banks of RA''}, published in CBA's monthly ``Statistical Overview''
  bulletin, broken down by source country. This is the series actually used: right frequency,
  and it is literally the input survey (``Report 24'') the official quarterly series is built
  from, so it is a defensible proxy even though it is not exactly the BPM6 definition (it
  includes some commercial transfers, not only personal remittances).
\end{enumerate}

\paragraph{Data collection.} 111 monthly bulletin issues (2015--01 through 2026--07) were
downloaded and converted to text; 69 parsed cleanly into a usable monthly series after two
rounds of chasing corrupted downloads and a site-side language mix-up (see below). Three
distinct URL-naming conventions were used by CBA's website across this period, and the
2023-onward English-language bulletins could not be located under any pattern tried -- some
requests silently returned a small placeholder file rather than a 404, and files repaired with
\texttt{qpdf} turned out on inspection to be the \emph{Armenian}-language edition, not corrupted
English ones. Coverage as obtained: complete 2015--2019, and complete January--September~2022
(the entire shock window); 2020 and most of 2021 are missing outright (bulletins never
successfully retrieved), and 2023 onward could not be recovered in English. Three spot checks
(2015--12, 2021--12, 2022--03) against the raw bulletin text confirm the parser is reading the
right numbers; a fourth check (the March-2021 value, read both from the March-2021 issue and
from the March-2022 issue's year-ago column) came back byte-identical, indicating this series is
not revised after publication -- unlike NBKR's.

\paragraph{Methodology.} A simplified real-time protocol
(\texttt{armenia\_realtime.py} / \texttt{armenia\_panel.py} / \texttt{armenia\_backtest.py}):
because
the series is not revised, ``known at issue date $v$'' only needs a fixed publication lag, not
NBKR's vintage/revision machinery. \textbf{The lag itself is an assumption, not a measured
fact}: one month, chosen because the bulletin for month $m$ already reports month $m$'s own
transfers, so the effective lag is however long the bulletin itself takes to compile -- no
official release calendar was found for this specific table (only for the quarterly series). Because 2020--2021 are missing, a seasonal (lag-12) feature is not usable for the shock-window
query rows that matter here (they would need the corresponding 2021 months); the Armenia model
therefore uses momentum lags only ($z_{t-1}, z_{t-2}, z_{t-3}$), a departure from the main
paper's $(1,2,3,12)$ specification. The comparison baseline is last-value naive (repeat the most
recently known month), not seasonal naive, for the same reason.

\paragraph{Result.} Same dev-tune-then-freeze protocol ($\alpha$ grid-searched on 2015--06
through 2019--12, frozen at $\alpha=0.3$, dev MAE \$14.5M).

\begin{table}[h]
\centering
\begin{tabular}{lc}
\toprule
 & Shock 2022 (n=8) \\
\midrule
Expanding ridge (lags 1--3) & 203.0 \\
Last-value naive & 163.4 \\
\bottomrule
\end{tabular}
\caption{MAE, million USD.}
\end{table}

\textbf{Naive beats the fitted model during the shock in Armenia too} -- the same qualitative
pattern as the Kyrgyzstan result in the main paper (a model extrapolating a recently-observed
trend is worst off exactly when that trend reverses). This is a real, if preliminary,
cross-country replication of the paper's central mechanism, obtained from an independent data
source, a different currency, a different country, and a different (undocumented, assumed)
publication-lag regime.

\subsection*{Caveats -- read before using this result for anything beyond a supervisor update}
\begin{itemize}
\item \textbf{$n=8$}, same small-sample caveats as the Kyrgyzstan shock window in the main paper,
  without that paper's bootstrap confidence intervals (not yet computed for Armenia).
\item \textbf{No ``ordinary period'' comparison exists.} The 2020--2021 gap sits exactly where a
  pre-shock evaluation window would go, so unlike the main paper's Table~1, there is no evidence
  yet on whether the fitted model beats naive \emph{outside} the shock for Armenia -- only that
  it loses \emph{during} it. Without that contrast, this result is suggestive of the mechanism
  but does not yet demonstrate the reversal the main paper documents for Kyrgyzstan.
\item \textbf{Publication lag is assumed, not measured} (see above).
\item \textbf{No seasonal term}, unlike the main paper's model.
\item \textbf{The series itself is a proxy} (total bank transfers, not strictly BPM6 personal
  remittances).
\item Coverage stops at January 2023; nothing is known about 2023--2026 Armenia behaviour.
\end{itemize}

\section{Where the original proposal's risks stand now}

The proposal itself flagged a risk: that the study would ``be perceived as an isolated case
study (anecdotal evidence)'' without a comparative panel. That risk is now partially addressed
-- there is a second, independent country showing the same qualitative pattern -- but not fully:
the Armenia result lacks the statistical rigor (bootstrap CIs, a proper main-window comparison,
a validated publication-lag assumption) that makes the Kyrgyzstan result in the MIDI paper
defensible on its own. It should be described, if used, as \emph{preliminary supporting
evidence}, not as a second confirmed case.

\section{Summary against the proposal}

\begin{table}[h]
\centering
\begin{tabular}{p{3.2cm}p{9.3cm}}
\toprule
Stage & Status \\
\midrule
1 (baseline) & Done -- MIDI 2026 submission, 8pp, reviewed. \\
2 (multimodal signals) & Done. Negative result: FX/Brent/IMOEX do not help, even hurt. \\
3a (Tajikistan) & Investigated, not feasible -- annual data only, and recent disclosure
  problems at the National Bank of Tajikistan. \\
3b (Armenia) & Data sourced (69 of 97 possible months, 2015--2023), pipeline built, first
  result obtained: same shock-period pattern as Kyrgyzstan. Preliminary -- see caveats above. \\
\bottomrule
\end{tabular}
\end{table}

\section{Recommended next steps, if this is pursued further}
\begin{enumerate}
\item Fill the Armenia 2020--2021 gap (needed for both a seasonal feature and a proper
  main-window comparison) -- requires a working internet connection to \texttt{cba.am}, which
  this project's sandboxed environment does not have; a further manual download round is needed
  the same way this round was done.
\item Confirm the assumed one-month Armenia publication lag against an actual CBA release
  calendar, if one can be found.
\item Compute bootstrap confidence intervals for the Armenia shock-window comparison, matching
  the main paper's treatment of the Kyrgyzstan result.
\item Decide, with the supervisor, whether Stage 2/3 material belongs in a revised/extended
  version of the MIDI paper (would require cutting other content to stay within any page limit)
  or is better positioned as a separate follow-up paper once Armenia's data and statistics are
  complete.
\end{enumerate}

\end{document}
